% Figure: Belleville (coned-disc) spring. Left, a section of the coned
% washer showing the free height h and the flattened state; right, the
% non-linear load-deflection curves for three height-to-thickness ratios,
% one rising monotonically, one nearly flat over a range (constant-force),
% and one with a snap-through dip (negative stiffness). The stress peaks
% at the inner top edge, the site the vignette sizes against.
\begin{figure}[htbp]
\centering
\begin{tikzpicture}
% Left: section sketch
\begin{scope}[shift={(0,0)}]
  % coned disc section (a shallow cone)
  \draw[thick, fill=engblue!6] (0,0) -- (1.6,0) -- (1.6,0.28) -- (0,1.1) -- (0,0.82) -- cycle;
  % free height marker
  \draw[engorange, {Latex}-{Latex}] (-0.3,0) -- (-0.3,1.1);
  \node[engorange, font=\scriptsize, anchor=east] at (-0.32,0.55) {$h$};
  % thickness marker
  \draw[engred, {Latex}-{Latex}] (1.75,0) -- (1.75,0.28);
  \node[engred, font=\scriptsize, anchor=west] at (1.78,0.14) {$t$};
  % load arrow at the top inner edge
  \draw[engblack, -{Latex}, thick] (0.1,1.5) -- (0.1,1.15);
  \node[font=\scriptsize, anchor=west] at (0.15,1.35) {$F$};
  % highest-stress corner marker
  \filldraw[engorange] (0,1.1) circle[radius=1.8pt];
  \node[engorange, font=\scriptsize, anchor=south west] at (0.02,1.12) {peak stress};
\end{scope}
% Right: load-deflection curves
\begin{scope}[shift={(3.6,0)}]
  \begin{axis}[
    width=8cm, height=6cm,
    xlabel={deflection $s/h$},
    ylabel={load $F$ (relative)},
    xmin=0, xmax=1.05, ymin=0, ymax=1.5,
    axis lines=left,
    tick label style={font=\scriptsize},
    label style={font=\small},
    legend pos=north west,
    legend style={font=\scriptsize},
  ]
  % nearly linear (h/t small):
  \addplot[engblue, very thick] coordinates {
    (0,0) (0.2,0.28) (0.4,0.55) (0.6,0.82) (0.8,1.08) (1.0,1.35)
  };
  \addlegendentry{$h/t \approx 0.4$ (stiff)}
  % constant-force plateau (h/t ~ 1.4):
  \addplot[engred, very thick] coordinates {
    (0,0) (0.2,0.55) (0.4,0.80) (0.6,0.85) (0.8,0.83) (1.0,0.95)
  };
  \addlegendentry{$h/t \approx 1.4$ (near-constant)}
  % snap-through, negative stiffness region (h/t ~ 2.8):
  \addplot[engorange, very thick] coordinates {
    (0,0) (0.15,0.70) (0.3,0.92) (0.45,0.80) (0.6,0.58) (0.75,0.62) (1.0,1.05)
  };
  \addlegendentry{$h/t \approx 2.8$ (snap-through)}
  \end{axis}
\end{scope}
\end{tikzpicture}
\caption[Belleville coned-disc spring]{A Belleville washer is a shallow
cone of disc that flattens under axial load, and its load-deflection
curve is set by the ratio of free height $h$ to thickness $t$. A low
ratio gives a nearly linear stiff spring; a ratio near $1.4$ gives a
near-constant force over a range of deflection, the property that holds a
bolt preload steady as a joint relaxes; a high ratio gives a
snap-through region of negative stiffness used in over-centre latches.
The peak stress is a bending stress at the inner top edge, marked in the
section, and it is that corner, not the average, that sizes the washer.}
\label{fig:vol03ch03:belleville-washer}
\end{figure}
